%-----------------------------------------------------------------------
%      This is not part of the manuscript itself
%-----------------------------------------------------------------------

%    Templates for common elements of a journal article;

%    Section headings
\section{}
\subsection{}

%    Ordinary theorem and proof
\begin{theorem}[Optional addition to theorem head]
% text of theorem
\end{theorem}

\begin{proof}[Optional replacement proof heading]
% text of proof
\end{proof}

%    Figure insertion; default placement is top; if the figure occupies
%    more than 75% of a page, the [p] option should be specified.
\begin{figure}
\includegraphics{filename}
\caption{text of caption}
\label{}
\end{figure}

% Numbered equation
\begin{equation}
\end{equation}

% Unnumbered equation
\begin{equation*}
\end{equation*}

% Aligned equations
\begin{align}
  &  \\
  &
\end{align}


% --- In case you use listings please use either of these examples as guidelines ---

% --- Example 1: Short Listing (No Page Breaks) ---
% For short code snippets where you want to prevent awkward page breaks,
% wrap the listing in a minipage. 
% Note: If the listing does not fit in the remaining space on the page, 
% the minipage will overflow rather than break. The 'float' option is a good alternative.
% Include a 'caption' if you want the listing to be numbered (e.g., "Listing 1") and referenced.
\noindent\begin{minipage}{\linewidth}
\begin{lstlisting}[caption={A short OpenFOAM snippet kept on a single page.}, label={lst:short_snippet}] 
// Calculate the flux
surfaceScalarField phi
(
    "phi",
    linearInterpolate(U) & mesh.Sf()
);
\end{lstlisting}
\end{minipage}
  
% --- Example 2: Long Listing (Allows Page Breaks) ---
% For longer code blocks that might span across multiple pages,
% DO NOT use a minipage. The listings package will handle the page
% breaks naturally. Again, a 'caption' is required to generate
% the "Listing X" numbering.
\begin{lstlisting}[caption={A longer OpenFOAM implementation that can safely break across pages.}, label={lst:piso_loop}]   
// Main PISO loop
for (int corr=0; corr<nCorr; corr++)
{
    volScalarField rAU(1.0/UEqn.A());
    volVectorField HbyA(constrainHbyA(rAU*UEqn.H(), U, p));
    
    surfaceScalarField phiHbyA
    (
        "phiHbyA",
        fvc::flux(HbyA)
      + fvc::interpolate(rAU)*fvc::ddtCorr(U, phi)
    );

    // Adjust pressure reference...
    adjustPhi(phiHbyA, U, p);
    
    // ... remainder of the loop ...
}
\end{lstlisting}